% Generated draft: requires mathematical and source review.
\documentclass[11pt]{article}
\usepackage[margin=25mm]{geometry}
\usepackage{fontspec}
\setmainfont{Latin Modern Roman}
\usepackage{amsmath,amssymb,mathtools,unicode-math}
\setmathfont{Latin Modern Math}
\usepackage{enumitem,needspace,xurl,hyperref}
\setlength{\parindent}{0pt}
\setlength{\parskip}{5pt}
\setlength{\emergencystretch}{3em}
\newcommand{\N}{\mathbb{N}}
\newcommand{\Z}{\mathbb{Z}}
\newcommand{\Q}{\mathbb{Q}}
\newcommand{\R}{\mathbb{R}}
\newcommand{\C}{\mathbb{C}}
\newcommand{\F}{\mathbb{F}}
\newcommand{\A}{\mathbb{A}}
\newcommand{\PP}{\mathbb{P}}
\newcommand{\E}{\mathbb{E}}
\newcommand{\eps}{\varepsilon}
\newcommand{\dd}{\,\mathrm{d}}
\begin{document}
\section*{UnsolvedMath ID: 30001666}
\textbf{Source problem code:} catalogue entry
\subsection{Definitions and mathematical statement}
An oriented tree is a finite directed graph obtained by assigning exactly one direction to each edge of a finite simple undirected tree. Interpret the supplied phrase ``directed tree'' in this sense, consistent with the supplied background's reference to oriented trees. For a tree with at least two vertices, interpret a leaf as a vertex of degree one in the underlying undirected tree, irrespective of edge directions. A tournament is a finite directed graph in which, for every two distinct vertices, exactly one of the two possible directed edges is present.

Under these interpretations, the nondegenerate formulation of the supplied question is: for every integer \(n\geq 2\), every oriented tree \(T\) with \(|V(T)|=n\), and
\[
k=\bigl|\{v\in V(T):\deg_{\mathrm{underlying}(T)}(v)=1\}\bigr|,
\]
does every tournament \(G\) with \(|V(G)|=n+k-1\) admit an injective map
\[
f:V(T)\longrightarrow V(G)
\]
such that
\[
(u,v)\in E(T)\quad\Longrightarrow\quad(f(u),f(v))\in E(G)
\qquad\text{for all }u,v\in V(T)?
\]
Here a copy means a subgraph copy preserving all edge directions; it need not be induced.

The supplied statement does not explicitly define leaves or impose \(n\geq2\). The displayed formulation states the case of trees with at least two vertices. Inclusion of the one-vertex tree requires clarification of the leaf convention and the resulting tournament order; that boundary case remains unresolved in this draft.
\subsection{Short English statement}
Must every tournament on \(n+k-1\) vertices contain every orientation of an \(n\)-vertex tree with \(k\) leaves, with all tree-edge directions preserved? Leaves are interpreted using underlying undirected degree. The source leaves the one-vertex boundary case unspecified.
\subsection{Sources}
\begin{itemize}
\item Local UnsolvedMath registry, supplied complete record for permanent ID \(30001666\), titled ``Directed Trees in Slightly Larger Tournaments,'' exact source code \texttt{catalogue entry}, in the Oberwolfach Reports Open Problems collection. The record assigns proposed year 2011 and supplies no individual proposer.
\item Supplied source citation: \emph{Combinatorics} (2011), Oberwolfach Reports, pp.\ 5--83, DOI \texttt{10.4171/owr/2011/01}. \url{https://doi.org/10.4171/owr/2011/01}. The supplied background describes these reports as documenting mathematical workshops and their open-problem discussions.
\item Supplied literature reference: ``Oberwolfach Report 01/2011,'' \url{https://doi.org/10.4171/owr/2011/01}. Literature checked at: \texttt{2026-08-21}. This is supplied dated curation, not an independently verified assessment. It labels the problem open, identifies the \(n+k-1\) universal bound for oriented trees with \(k\) leaves as Conjecture 3, and reports that only special cases were known. The supplied background explicitly requests independent verification of time-sensitive status and cited claims.
\end{itemize}
\end{document}
